\practicechapter{第 13 章实践：System V AMD64 ABI}

本章对应第一册第 13 章“System V AMD64 ABI”。ABI 是 \term{Application Binary Interface}{应用二进制接口}，它规定二进制层面的函数调用、寄存器使用、栈对齐、符号和链接约定。本章不要求一次掌握全部 ABI，而是从符号表和一个未内联函数开始，把源码函数连接到二进制可见边界。

\practicesection{一、对应原书问题：ABI 不只是在说寄存器}

很多人学 ABI 时只记住“前几个整数参数放在 \code{rdi}、\code{rsi}、\code{rdx}...”。这当然重要，但 ABI 还体现在符号、调用边界、栈、返回值、链接器可见性等地方。实践卷的入口是 symbol table，因为它能告诉你某个函数或对象是否在当前二进制中作为符号出现。

\code{SymbolSummary} 记录 name、section name、type、binding、section index、value 和 size。它不是完整 ABI 模型，但足以回答几个关键问题：这个名字是否可见，它属于哪个 section，它是函数还是对象，它的地址值是什么，它是否可能被优化、strip 或 name mangling 改变。

\practicesection{二、从特例推到规律：源码函数不一定有同名符号}

特例一：C++ 函数 \code{int add_i32(int,int)} 可能在符号表中出现为经过 name mangling 的名字，而不是源代码中的简单名字。特例二：函数被 inline 后，最终可执行文件中可能没有独立函数符号。特例三：用 \code{extern "C"} 声明的函数名字通常更接近源码名字，但仍受可见性、链接和优化影响。

由此推出规律：ABI 观察要同时看源码声明、编译命令、符号表、反汇编和优化等级。符号表不能单独证明所有调用约定，但它能告诉我们当前二进制暴露了哪些边界。想观察寄存器传参，还需要反汇编或调试器。

\practicesection{三、实践任务：保留一个 ABI 观察点}

写一个不容易被内联的小函数：

\begin{lstlisting}[language=C++]
extern "C" __attribute__((noinline))
int abi_add_i32(int a, int b) {
    return a + b;
}
\end{lstlisting}

编译后观察：

\begin{lstlisting}[language=bash]
g++ -std=c++20 -O2 -g abi_sample.cpp -o abi-sample
cpu1ee=books/cpu-volume-1-practice/build/executable-evidence-debug/labs/executable_evidence/cpu1ee_cli
$cpu1ee ./abi-sample > reports/abi-sample.report
readelf -s ./abi-sample > reports/abi-sample.symbols.txt
objdump -d -Mintel ./abi-sample > reports/abi-sample.disasm.txt
\end{lstlisting}

在 \cmd{readelf -s} 中找 \code{abi\_add\_i32}，记录 symbol value、size、type、binding、section index。再在反汇编中找函数入口，观察参数是否从约定寄存器进入计算。若编译器仍然改变形态，报告要写实际观察结果，而不是强行套模板。

\practicesection{四、验收：ABI 结论要有层级}

本章交付物是 \filepath{reports/ch13-system-v-abi.md}。通过标准如下：

\begin{enumerate}
  \item 能解释 ABI 全称和它解决的问题。
  \item 报告中至少有一个未内联函数的 symbol 证据。
  \item 能说明 type、binding、section index、value、size 的含义。
  \item 能用反汇编观察至少一个参数寄存器或返回寄存器。
  \item 能说明本卷解析器没有完整验证调用约定，它只是提供符号和 section 证据的一部分。
\end{enumerate}

这章完成后，读者应该把 ABI 看成一组可观察的二进制约定，而不是一张寄存器背诵表。
